% The Collected Mathematical Papers of Arthur Cayley, Vol. XI
% Transcription of PDF pages 426--437 (printed pages 402--413)
% Papers covered:
%   777  A SOLVABLE CASE OF THE QUINTIC EQUATION                                pp.402--404
%   778  [ADDITION TO MR HUDSON'S PAPER "ON EQUAL ROOTS OF EQUATIONS."]         pp.405--407
%   779  [NOTE ON MR JEFFERY'S PAPER ON CERTAIN QUARTIC CURVES...]              p.   408
%   780  [ADDITION TO MR HAMMOND'S PAPER "NOTE ON AN EXCEPTIONAL CASE..."]      pp.409--410
%   781  ON THE AUTOMORPHIC TRANSFORMATION OF THE BINARY CUBIC FUNCTION  (part) pp.411--413
% Scan-verified against collectedmathema11cayluoft.pdf at 200 dpi.

\documentclass[11pt]{article}
\usepackage[a4paper,margin=0.65in]{geometry}
\usepackage[T1]{fontenc}
\usepackage{amsmath,amssymb,amsfonts}
\usepackage{graphicx}
\usepackage{pdfpages}
\usepackage[english,shorthands=off]{babel}
\usepackage{fancyhdr}
\usepackage[bookmarks=false,colorlinks=false]{hyperref}
\usepackage[protrusion=true,expansion=false]{microtype}
\usepackage{array}
\usepackage{longtable}
\emergencystretch=3em
\tolerance=600
\providecommand{\modd}{\mathrm{mod}}
\providecommand{\Disc}{\mathrm{Disc}}
\providecommand{\canont}{}
\providecommand{\asterism}{*}
\providecommand{\Hessian}{H}
\providecommand{\tome}{}
\providecommand{\page}{p.}
\pagestyle{fancy}\fancyhf{}%
\fancyhead[L]{Pages 426--437}\fancyhead[R]{Cayley --- Collected Papers, Vol.\ XI}%
\renewcommand{\headrulewidth}{0.4pt}
\setlength{\headheight}{15pt}

\allowdisplaybreaks

\begin{document}

% ============================================================
% Paper 777.  A SOLVABLE CASE OF THE QUINTIC EQUATION.
% PDF pp.\ 426--428 (printed pp.\ 402--404).
% ============================================================

\begin{center}
\textbf{777.}\\[4pt]
\textsc{A SOLVABLE CASE OF THE QUINTIC EQUATION.}\\[4pt]
[From the \textit{Quarterly Journal of Pure and Applied Mathematics}, vol.\ \textsc{xviii}.\ (1882), pp.\ 154--157.]
\end{center}

\bigskip

% ----- PDF p426 (printed 402) -----
\noindent
The roots of the general quintic equation
\[
  (a,\,b,\,c,\,d,\,e,\,f \mid x,\,1)^{5} = 0
\]
may be taken to be
\begin{align*}
  & -\tfrac{b}{a} + B + C + D + E,\\
  & -\tfrac{b}{a} + \omega B + \omega^{2} C + \omega^{3} D + \omega^{4} E,\\
  & -\tfrac{b}{a} + \omega^{2} B + \omega^{4} C + \omega D + \omega^{3} E,\\
  & -\tfrac{b}{a} + \omega^{3} B + \omega C + \omega^{4} D + \omega^{2} E,\\
  & -\tfrac{b}{a} + \omega^{4} B + \omega^{3} C + \omega^{2} D + \omega E,
\end{align*}
where $\omega$ is an imaginary fifth root of unity; and if one of the four
functions $B,\,C,\,D,\,E$ is $=0$, say if $E = 0$ (this implies of course a
single relation between the coefficients), then the equation is solvable.

\medskip
Writing $x = \xi - \dfrac{b}{a}$, we have
\[
  (a,\,b,\,c,\,d,\,e,\,f \mid \xi - \tfrac{b}{a},\,1)^{5}
   = (a',\,0,\,c',\,d',\,e',\,f' \mid \xi,\,1)^{5},
\]
where
\begin{align*}
  a' &= a,\\
  ac' &= ac - b^{2},\\
  a^{2} d' &= a^{2} d - 3abc + 2b^{3},\\
  a^{3} e' &= a^{3} e - 4a^{2} b d + 6 a b^{2} c - 3 b^{4},\\
  a^{4} f' &= a^{4} f - 5 a^{3} b e + 10 a^{2} b^{2} d - 10 a b^{3} c + 4 b^{5},
\end{align*}
and the roots of the new equation
\[
  (a',\,0,\,c',\,d',\,e',\,f' \mid \xi,\,1)^{5} = 0
\]

% ----- PDF p427 (printed 403) -----
have the above-mentioned values, omitting therefrom the term $-\dfrac{b}{a}$:
we find without difficulty
\begin{align*}
  2\frac{c'}{a} &= - BE - CD,\\
  2\frac{d'}{a} &= - B^{2} D - BC^{2} - CE^{2} - D^{2} E,\\
  \frac{e'}{a} &= - B^{3} C - B^{2} E^{2} + BCDE + BD^{3} + C^{3} E + C^{2} D^{2} - DE^{3},\\
  \frac{f'}{a} &= - B^{5} + 5 B^{2} D E - 5 B^{2} C^{2} E - 5 B^{2} C D^{2} + 5 B C^{3} D + 5 B C E^{3}\\
        &\qquad - 5 B D^{3} E^{2} - C^{5} + 5 C D^{3} E - 5 C D^{2} E^{2} - D^{5} - E^{5},
\end{align*}
and hence, when $E = 0$, we have
\begin{align*}
  2\frac{c'}{a} &= - CD,\\
  2\frac{d'}{a} &= - B^{2} D - BC^{2},\\
  \frac{e'}{a} &= - B^{3} C - BD^{3} - C^{2} D^{2},\\
  \frac{f'}{a} &= - B^{5} - 5 B^{2} C D^{2} + 5 B C^{3} D - C^{5} - D^{5},
\end{align*}
or, as these may be written,
\begin{align*}
  - 2 \tfrac{c'}{a} &= CD,\\
  - 2 \tfrac{d'}{a} &= B^{2} D + BC^{2},\\
  - \tfrac{e'}{a} - \tfrac{c'^{2}}{a^{2}} &= B^{3} C - BD^{3},\\
  - \tfrac{f'}{a} &= B^{5} + C^{5} + D^{5} - 10\,\tfrac{c'}{a}\,(B^{2} D - BC^{2}),
\end{align*}
equations which imply a single relation between the coefficients $a',\,c',\,d',\,e',\,f'$.
Supposing this satisfied, we may attend only to the first three equations; or, writing
for convenience,
\begin{align*}
  \gamma &= - 2\tfrac{c'}{a} = -\tfrac{2}{a^{2}}(ac - b^{2}),\\
  \delta &= - 2\tfrac{d'}{a} = -\tfrac{2}{a^{3}}(a^{2} d - 3 a b c + 2 b^{3}),\\
  \theta &= -\tfrac{e'}{a} - \tfrac{c'^{2}}{a^{2}}
          = -\tfrac{1}{a^{4}}\bigl\{a^{2}(ae - 4 b d + 3 c^{2}) + (ac - b^{2})^{2}\bigr\},
\end{align*}
the equations are
\begin{align*}
  \gamma &= CD,\\
  \delta &= B(BD + C^{2}),\\
  \theta &= B(D^{3} - B^{2} C).
\end{align*}
\hfill $51{-}2$

% ----- PDF p428 (printed 404) -----
The first equation gives $C = \dfrac{\gamma}{D}$, and substituting this value in the other two
equations, we have
\begin{align*}
  B^{2} D^{2} + B \gamma^{2} - \delta D &= 0,\\
  B^{3} \gamma^{2} + B D^{4} + \theta D &= 0.
\end{align*}

Eliminating $B$, the result is obtained in the form $\det. = 0$, where in the last
column of the determinant each term is divisible by $D$; and omitting this factor,
the result is
\[
  \begin{vmatrix}
    D^{2}, & \gamma^{2}, & 0, & -\delta D & \\
    0, & D^{2}, & \gamma^{2}, & -\delta D & \\
    \gamma^{2}, & 0, & D^{4}, & \theta & \\
    0, & \gamma^{2}, & 0, & D^{4}, & \theta \\
  \end{vmatrix} = 0.
\]

If, in order to develop the determinant, we consider it as a sum of products,
each first factor being a minor composed out of columns 1 and 2, and the second
factor being the complementary minor composed out of columns 3, 4, 5 (the several
products being of course taken each with its proper sign), the expansion presents
itself in the form
\begin{align*}
  & D^{4}(-\gamma^{2} \theta D^{2} + \theta^{2} D^{3})\\
  & {} - D^{2} \gamma^{2}(- \theta \gamma^{2} D^{2} + \theta D^{5} - \delta^{2} D^{3})\\
  & {} - \gamma^{2} D^{2}(- \delta D^{6} - \theta \gamma^{4})\\
  & {} + \gamma^{4}(- \gamma^{4} \delta D^{2} - \delta^{3} D^{4} - \delta \gamma^{2} D^{4})\\
  & {} - \gamma^{8} D^{2}.
\end{align*}

Hence, collecting, and changing the sign of the whole expression, we obtain
\[
  B D^{6} - (2 \gamma^{2} \delta + \theta^{2}) D^{4} + (-\gamma^{4} \delta + \delta^{2}\gamma^{2}\theta + \delta^{2}) D^{3} + \gamma^{8} = 0,
\]
\emph{viz.} a cubic equation for $D^{2}$.
We have then as above $C = \dfrac{\gamma}{D}$, and $B$ is given rationally
as the common root of the foregoing quadric and cubic equations satisfied by $B$.

Substituting for $\gamma,\,\delta,\,\theta$ their values in terms of the original coefficients, the
equation for $D^{2}$ becomes
\[
  \resizebox{\textwidth}{!}{$\displaystyle
  \left\{
  \begin{aligned}
    & 2(a^{2} d - 3 a b c + 2 b^{3})(aD)^{6}\\
    & {} + a^{2}\bigl[(ae - 4 b d + 3 c^{2})^{2} a^{2} + a^{2}(ac - b^{2})^{2}(ae - 4 b d + 3 c^{2}) - 16(ac - b^{2})^{3}(a^{2} d - 3 a b c + 2 b^{3})\bigr] (aD)^{4}\\
    & {} + \bigl[28(ac - b^{2})^{2}(a^{2} d - 3 a b c + 2 b^{3})\\
    & \qquad {} + 4(ae - b^{2})^{2}\bigl\{4 a^{2}(ac - b^{2})(a^{2} d - 3 a b c + 2 b^{3})(ae - 4 b d + 3 c^{2}) + 8(a^{2} d - 3 a b c + 2 b^{3})^{2}\bigr\}\bigr] (aD)^{3}\\
    & {} - 128 (ac - b^{2})^{3} \bigl\{a^{2}(ae - 4 b d + 3 c^{2})(ac^{2} - b^{2})^{2}\bigr\} = 0,
  \end{aligned}
  \right.
  $}
\]
and the solution of the given quintic equation thus ultimately depends upon that of
this cubic equation.

% ============================================================
% Paper 778.  [ADDITION TO MR HUDSON'S PAPER "ON EQUAL ROOTS OF EQUATIONS."]
% PDF pp.\ 429--431 (printed pp.\ 405--407).
% ============================================================

\bigskip

\begin{center}
\textbf{778.}\\[4pt]
\textsc{[ADDITION TO MR HUDSON'S PAPER ``ON EQUAL ROOTS OF EQUATIONS.'']}\\[4pt]
[From the \textit{Quarterly Journal of Pure and Applied Mathematics}, vol.\ \textsc{xviii}.\ (1882), pp.\ 226--229.]
\end{center}

\bigskip

% ----- PDF p429 (printed 405) -----
\noindent
It seems desirable to present in a more developed form some of the results of
the foregoing paper.

\medskip
Thus, if the equation $(a_{0},\,a_{1},\,\ldots,\,a_{n} \mid x,\,1)^{n} = 0$
of the order $n$ has $n - v$ equal roots, where $v$ is not $> \tfrac{1}{2}n - 1$, then we have
$\psi(r,\,v+1,\,m) = 0$, where $m$ has any one of the values $0,\,1,\,\ldots,\,n - 2v - 2$,
and $r$ any one of the values
\[
  2v + 2,\; 2v + 3,\; \ldots,\; n - m.
\]

The signification is
\begin{align*}
  \psi(r,\,v+1,\,m) &= r \cdot 1 \cdot \frac{1}{[r]^{v+1}} a_{m} a_{r+m}\\
  &\quad - (r - 2) \cdot \frac{v+1}{1} \cdot \frac{1}{[r-1]^{v+1}} a_{m+1} a_{r+m-1}\\
  &\quad + (r - 4) \cdot \frac{v+1\,.\,v+2}{1\,.\,2} \cdot \frac{1}{[r-2]^{v+1}} a_{m+2} a_{r+m-2}\\
  &\quad + (-)^{s} (r - 2s) \cdot \frac{[v+1]^{s}}{[s]^{s}} \cdot \frac{1}{[r-s]^{v+1}} a_{m+s} a_{r+m-s}\\
  &\quad \cdots\\
  &\quad + (-)^{v+1}(r - 2v - 2) \cdot \frac{1}{[r-v-1]^{v+1}} a_{m+v+1} a_{r+m-v-1}.
\end{align*}

Thus, when $v = 0$, the condition is
\[
  r \cdot \frac{1}{r \cdot r - 1} a_{m} a_{r+m} - (r - 2) \cdot \frac{1}{r - 1 \cdot r - 2} a_{m+1} a_{r+m-1} = 0,
\]

% ----- PDF p430 (printed 406) -----
that is,
\[
  a_{m} a_{r+m} - a_{m+1} a_{r+m-1} = 0,
\]
satisfied when the equation has all its roots equal.

\medskip
The values of $m$ are $0,\,1,\,2,\,\ldots,\,n - 2$, and those of $r$ are
$2,\,3,\,\ldots,\,n - m$; in particular, if $m = 0$, the values of $r$ are
$2,\,3,\,\ldots,\,n$, and the corresponding conditions are
\begin{align*}
  a_{0} a_{2} - a_{1}^{2} &= 0,\\
  a_{0} a_{3} - a_{1} a_{2} &= 0,\\
  a_{0} a_{n} - a_{1} a_{n-1} &= 0,
\end{align*}
and so for the different values of $m$ up to the final value $r = 2$, for which $m = n - 2$,
and the condition is
\[
  a_{n-2} a_{n} - a_{n-1}^{2} = 0;
\]
we have thus, it is clear, the whole series of conditions included in
\[
  \begin{vmatrix}
    a_{0}, & a_{1}, & a_{2}, & \ldots, & a_{n-1}\\
    a_{1}, & a_{2}, & a_{3}, & \ldots, & a_{n}
  \end{vmatrix} = 0,
\]
which are obviously satisfied in the case in question of the roots being all equal.

\medskip
Again, when $v = 1$, the condition for $n - 1$ equal roots is
\[
  \resizebox{\textwidth}{!}{$\displaystyle
  r \cdot 1 \cdot \frac{1}{r \cdot r - 1 \cdot r - 2} a_{m} a_{r+m}
  - (r - 2) \cdot \frac{2}{1} \cdot \frac{1}{r - 1 \cdot r - 2 \cdot r - 3} a_{m+1} a_{r+m-1}
  + (r - 4) \cdot 1 \cdot \frac{1}{r - 2 \cdot r - 3 \cdot r - 4} a_{m+2} a_{r+m-2} = 0;
  $}
\]
that is,
\[
  r - 1 \cdot r - 2 - r - 1 \cdot r - 3 + r - 2 \cdot r - 3 = 0;
\]
or, what is the same thing,
\[
  (r - 3)\,a_{m} a_{r+m} - 2(r - 2)\,a_{m+1} a_{r+m-1} + (r - 1)\,a_{m+2} a_{r+m-2} = 0,
\]
where $n = 4$ at least, and $m,\,r$ have the values
\[
  m = 0,\,1,\,2,\,\ldots,\,n - 4
\]
\[
  r = \begin{cases} 4,\,4,\,\ldots,\,4\\ 5,\,5\\ \vdots\\ n - 1\\ n\end{cases}
\]
thus, when $n = 4$, the only values are $m = 0$, $r = 4$, and the condition is
\[
  a_{0} a_{4} - 4 a_{1} a_{3} + 3 a_{2}^{2} = 0.
\]

% ----- PDF p431 (printed 407) -----
Similarly, when $v = 2$, the condition for $n - 2$ equal roots is found to be
\[
  \resizebox{\textwidth}{!}{$\displaystyle
  \frac{a_{m} a_{r+m}}{r - 1 \cdot r - 2 \cdot r - 3}
  - \frac{3 a_{m+1} a_{r+m-1}}{r - 1 \cdot r - 2 \cdot r - 4}
  + \frac{3 a_{m+2} a_{r+m-2}}{r - 2 \cdot r - 3 \cdot r - 5}
  - \frac{a_{m+3} a_{r+m-3}}{r - 3 \cdot r - 4 \cdot r - 5} = 0;
  $}
\]
or, what is the same thing,
\begin{align*}
  & r - 4 \cdot r - 5 \cdot a_{m} a_{r+m}\\
  & {} - 3 \cdot r - 2 \cdot r - 5 \cdot a_{m+1} a_{r+m-1}\\
  & {} + 3 \cdot r - 1 \cdot r - 4 \cdot a_{m+2} a_{r+m-2}\\
  & {} - r - 1 \cdot r - 2 \cdot a_{m+3} a_{r+m-3} = 0,
\end{align*}
where $n = 6$ at least, and $m,\,r$ have the values
\[
  m = 0,\,1,\,\ldots,\,n - 6
\]
\[
  r = \begin{cases} 6,\,6,\,\ldots,\,6\\ 7,\,7\\ \vdots\\ n - 1\\ n\end{cases}
\]

Observe that the sum of the coefficients is $0$, \emph{viz.}
\[
  (r - 4)(r - 5) - 3(r - 2)(r - 5) + 3(r - 1)(r - 4) - (r - 1)(r - 2) = 0,
\]
this should obviously be the case, since the conditions for $n - 2$ equal roots must
be satisfied when the roots are all of them equal; and the property serves as a
verification.

\medskip
It is to be remarked that the equation $\psi(r,\,v + 1,\,m) = 0$ does not in all cases
give all the conditions for the existence of $n - v$ equal roots in an equation of the
order $n$; thus when $n = 3$ and $v = 1$, we cannot by means of it obtain the condition that
a cubic equation may have $2$ equal roots. The problem really considered is that of
the determination of those quadric functions of the coefficients which vanish in the
case of $n - v$ equal roots; and in the case in question $(n = 3,\,v = 1)$ there is no
quadric function which vanishes, but the condition depends on a cubic function.

The question of the quadric functions which vanish in the case of $n - v$ equal
roots, and to a small extent that of the cubic functions which thus vanish, is considered
in Dr Salmon's ``Note on the conditions that an equation may have equal roots,''
\textit{Camb.\ and Dublin Math.\ Jour.}, t.\ \textsc{v}.\ (1850), pp.\ 159--165, and in particular the
equation there obtained p.\ 161 is the equation $\psi(0,\,v + 1,\,n) = 0$.

% ============================================================
% Paper 779.  [NOTE ON MR JEFFERY'S PAPER...]
% PDF p.\ 432 (printed p.\ 408).
% ============================================================

\bigskip

\begin{center}
\textbf{779.}\\[4pt]
\textsc{[NOTE ON MR JEFFERY'S PAPER ``ON CERTAIN QUARTIC CURVES WHICH HAVE A CUSP AT INFINITY, WHEREAT THE LINE AT INFINITY IS A TANGENT.'']}\\[4pt]
[From the \textit{Proceedings of the London Mathematical Society}, vol.\ \textsc{xiv}.\ (1883), p.\ 85.]
\end{center}

\bigskip

% ----- PDF p432 (printed 408) -----
\noindent
The assumed form $\kappa \phi^{2} = u_{4}$, or, as this is afterwards written,
\[
  2 y^{3} y = a x^{2} + 2 b x y + c y^{2} + 2 e x + 2 d y + \lambda,
\]
is, I think, introduced without a proper explanation. Say, the form is
$x^{2} y = z^{2}\,(*\!\mid\!x,\,y,\,z)^{2}$, it ought to be shown how for a cuspidal quartic
we arrive at this form; \emph{viz.} taking the cusp to be at the point $(x = 0,\,z = 0)$,
$z = 0$ for the tangent at the cusp, and $x = 0$ an arbitrary line through the cusp;
then the line $z = 0$ besides intersects the curve in a single point, and, if $y = 0$
is taken as the tangent at that point, the equation of the curve must, it can be
seen, be of the form
\[
  (ax^{2} + exz) y = z^{2}\,(a,\,b,\,c,\,f,\,g,\,h \mid x,\,y,\,z)^{2}.
\]
The conic $(a,\,b,\,c,\,f,\,g,\,h \mid x,\,y,\,z)^{2}$ touches the quartic at each of the two
intersections of the quartic with the arbitrary line $x = 0$; and we cannot, so long
as the line remains arbitrary, find a conic which shall osculate the quartic at the
two points in question; but, for the particular line $x + \tfrac{1}{2} d z = 0$, there
exists such a conic, \emph{viz.} writing $x$ instead of $x + \tfrac{1}{2} d z$, the form
is $x^{2} y = z^{2}\,(a',\,b',\,c',\,f',\,g',\,h' \mid x,\,y,\,z)^{2}$, and the new conic
$(a',\,\ldots \mid x,\,y,\,z)^{2} = 0$ has the property in question. This is the adopted form,
and it thus appears that in it the line $x = 0$ is a determinate line, \emph{viz.} the line
passing through the cusp and the two points of osculation of the osculating conic.

It thus appears that in the assumed form the lines $x = 0,\,y = 0,\,z = 0$ are determinate
lines.

% ============================================================
% Paper 780.  [ADDITION TO MR HAMMOND'S PAPER...]
% PDF pp.\ 433--434 (printed pp.\ 409--410).
% ============================================================

\bigskip

\begin{center}
\textbf{780.}\\[4pt]
\textsc{[ADDITION TO MR HAMMOND'S PAPER ``NOTE ON AN EXCEPTIONAL CASE IN WHICH THE FUNDAMENTAL POSTULATE OF PROFESSOR SYLVESTER'S THEORY OF TAMISAGE FAILS.'']}\\[4pt]
[From the \textit{Proceedings of the London Mathematical Society}, vol.\ \textsc{xiv}.\ (1883), pp.\ 88--91. Read Dec.\ 14, 1882.]
\end{center}

\bigskip

% ----- PDF p433 (printed 409) -----
\noindent
The extreme importance of Mr Hammond's result, as regards the entire subject
of Covariants, leads me to reproduce his investigation in the notation of my Memoirs
on Quantics, and with a somewhat different arrangement of the formul\ae. For the
binary seventhic
\[
  (a,\,b,\,c,\,d,\,e,\,f,\,g,\,h \mid x,\,y)^{7},
\]
the four composite seminvariants of the deg-weight $5\,.\,11$ (sources of covariants of the
deg-order $5\,.\,13$) are

\medskip
\begin{center}
\begin{tabular}{|c|c|c|c|}
\hline
\textsc{I}. & \textsc{II}. & \textsc{III}. & \textsc{IV}.\\
\hline
$1\,.\,7$ & $4\,.\,6$ & $2\,.\,10$ & $3\,.\,3$ \quad Deg-order.\\
$1\,.\,0$ & $4\,.\,11$ & $2\,.\,2$ & $3\,.\,9$ \quad Deg-weight.\\
\hline
$a + 1$ &
$\begin{array}{l} a^{2} e h + 1\\ a^{2} f g - 1\\ a b d h - 4\\ a b e g - 2\\ a b f^{2} + 6\\ a c^{2} h - 3\\ a c d g - 2\\ a c e f - 6\\ a d^{2} f + 10\\ a d^{2} e - 5\\ a d e^{2} + \ldots\\ b^{2} d g + 20\\ b^{2} e f + 57\\ b c^{2} g - 15\\ b c d f - 24\\ b c e^{2} - 30\\ b c^{2} e - 10\\ c^{4} f + 27\\ c^{3} d e - 45\\ c d^{3} + 20\end{array}$ &
$\begin{array}{l} ac + 1\\ b^{2} - 1\end{array}$ &
$\begin{array}{l} ach + 2\\ adg - 7\\ aef + 5\\ ac^{2}fh - 2\\ beg + 7\\ bdf + 22\\ be^{2} - 25\\ cf - 27\\ cde + 45\\ d^{3} - 20\end{array}$\\
\hline
\end{tabular}
\end{center}

\bigskip

% ----- PDF p434 (printed 410) -----
\noindent
\begin{center}
\begin{tabular}{|c|c|c|}
\hline
\textsc{II}. (continued) & \textsc{III}. (continued) & \textsc{IV}. (continued)\\
\hline
$2\,.\,6$ & $3\,.\,7$ & $3\,.\,11$ \quad Deg-order.\\
$2\,.\,4$ & $3\,.\,1$ & $3\,.\,5$ \quad Deg-weight.\\
\hline
$\begin{array}{l} ae + 1\\ bd - 4\\ c^{2} + 3\end{array}$ &
$\begin{array}{l} a^{2} h + 1\\ abg - 7\\ acf + 9\\ ade - 5\\ b^{2} f + 12\\ bce - 30\\ bd^{2} + 20\end{array}$ &
$\begin{array}{l} ag + 1\\ bf - 6\\ ce + 16\\ d^{2} - 10\end{array}$\\
\hline
$\begin{array}{l} a^{2} f + 1\\ abe - 5\\ ad - 2\\ ab^{2} d - 3\\ be^{2} - 6\end{array}$ & & \\
\hline
\end{tabular}
\end{center}

\medskip
and it is here at once obvious that there exists a syzygy of the form $\text{I} = \text{III} - \text{IV}$;
in fact, if in \textsc{III}.\ and \textsc{IV}.\ we write $a = 0$, then the values are each
\[
  = -2b\,(2 b d - 3 c^{2})\,(6 b f - 15 c e + 10 d^{2})
\]
hence $\text{III} - \text{IV}$ must divide by $a$, the quotient being a seminvariant of the deg-weight
$4\,.\,11$, which can only be a numerical multiple of the second factor of \textsc{I}., and is in
fact this second factor, that is, we have the syzygy $\text{I} = \text{III} - \text{IV}$.

\medskip
Working out the values of the four products, and joining to them the expression
for the irreducible seminvariant of the same deg-weight $5\,.\,11$ ($\Omega,\,\Theta^{*}$ of my tables [774]
for the binary sextic), we have the table:

\medskip
\begin{center}
\begin{tabular}{|l|l|r|r|r|r|}
\hline
$5\,.\,10$ & $5\,.\,11$ & I. & III. & IV. & II.\\
\hline
$a^{2}dh$ & $a^{2}efh$ & $+1$ & $+1$ & & \\
$a^{2}eg$ & $a^{2}fg$ & $-1$ & $-4$ & $+1$ & \\
$a^{2}f^{2}$ & $a^{2}bdh$ & $-4$ & $-7$ & & \\
& $a^{2}beg$ & $-2$ & $+3$ & $-5$ & \\
$a^{2}bch$ & $a^{2}bf^{2}$ & $+3$ & $+9$ & $-6$ & \\
$a^{2}bdg$ & $a^{2}c^{2}h$ & $-2$ & $-5$ & $+2$ & \\
$a^{2}bef$ & $a^{2}cdg$ & $-1$ & $+28$ & $+15$ & $+2$\\
$a^{2}c^{2}g$ & $a^{2}cef$ & $+10$ & $+12$ & $-10$ & $-7$\\
$a^{2}cdf$ & $a^{2}d^{2}f$ & $+3$ & & & $+5$\\
$a^{2}ce^{2}$ & $a^{2}de^{2}$ & $-5$ & $-21$ & & $-4$\\
$a^{2}d^{2}e$ & $ab^{2}ch$ & $-2$ & $-36$ & & \\
$ab^{3}h$ & $ab^{2}cg$ & $+20$ & $-30$ & $+8$ & $+7$\\
$ab^{2}cg$ & $ab^{2}df$ & $+1$ & $+67$ & $-45$ & $-6$\\
$ab^{2}df$ & $ab^{2}e^{2}$ & $-15$ & $+40$ & & \\
$ab^{2}e^{2}$ & $abc^{2}g$ & $-24$ & $+27$ & $-6$ & $+7$\\
$abc^{2}f$ & $abc^{2}f$ & $-30$ & $-15$ & $-12$ & $+22$\\
$abcde$ & $abcdf$ & $+27$ & $-48$ & & \\
$abd^{3}$ & $abce^{2}$ & $-14$ & $-45$ & & $-25$\\
$ac^{3}e$ & $abd^{2}e$ & $+20$ & $+36$ & $+50$ & $-27$\\
$ac^{2}d^{2}$ & $ac^{3}f$ & $+11$ & $+120$ & $+30$ & $+46$\\
$a^{2}c^{2}g$ & $ac^{2}de$ & $+1$ & $-80$ & $-20$ & $-20$\\
$ab^{2}cf$ & $ac d^{3}$ & $+9$ & $-90$ & & $-2$\\
$ab^{2}de$ & $a^{2}b^{2}h$ & $-14$ & $+60$ & $-48$ & $-7$\\
$ab^{2}c^{2}e$ & $b^{3}cg$ & $+6$ & & $-22$ & $+26$\\
$ab^{2}cd^{2}$ & $b^{3}df$ & $+8$ & & $+86$ & $+27$\\
$abc^{3}d$ & $b^{3}e^{2}$ & $-9$ & & $+120$ & $-45$\\
& $b^{2}c^{2}f$ & $-6$ & & $-80$ & $+20$\\
& $b^{2}cde$ & $+16$ & & $-90$ & \\
& $b^{2}d^{3}$ & $-8$ & & $+60$ & \\
& $bc^{3}e$ & $-3$ & & & \\
& $bc^{2}d^{2}$ & $+2$ & & & \\
& $c^{4}d$ & & & & \\
\hline
\end{tabular}
\end{center}

\medskip
I have prefixed to the table the literal terms of the deg-weight $5\,.\,10$; for the
deg-weights $5\,.\,11$ and $5\,.\,10$, the numbers of terms are $= 30$ and $26$ respectively;
and it is the difference of these $30 - 26 = 4$, which gives the number of asyzygetic
seminvariants of the deg-weight $5\,.\,11$.

% ============================================================
% Paper 781.  ON THE AUTOMORPHIC TRANSFORMATION OF THE BINARY CUBIC FUNCTION.
% PDF pp.\ 435--437 (printed pp.\ 411--413).
% ============================================================

\bigskip

\begin{center}
\textbf{781.}\\[4pt]
\textsc{ON THE AUTOMORPHIC TRANSFORMATION OF THE BINARY CUBIC FUNCTION.}\\[4pt]
[From the \textit{Proceedings of the London Mathematical Society}, vol.\ \textsc{xiv}.\ (1883), pp.\ 103--108. Read Jan.\ 11, 1883.]
\end{center}

\bigskip

% ----- PDF p435 (printed 411) -----
\noindent
I CONSIDER the cubic equation $(a,\,b,\,c,\,d \mid x,\,1)^{3} = 0$. It is shown (Serret,
\textit{Cours d'Alg\`ebre sup\'erieure}, 4th ed., Paris, 1879, t.\ II.\ pp.\ 466--471) how, given one
root of the equation, the other two roots can be each of them expressed rationally
in terms of this root and of the square root of the discriminant; \emph{viz.} making
the proper changes of notation, and writing
\begin{gather*}
  A,\,B,\,C = ac - b^{2},\;\; ad - bc,\;\; bd - c^{2},\;\; \lambda = \sqrt{-1},\\
  \Omega = B^{2} - 4AC,\quad
   = a^{2} d^{2} - 4 a c^{3} + 6 a b c d - 4 b^{3} d - 3 b^{2} c^{2},\\
  \alpha = \frac{-\lambda \sqrt{\Omega} + B}{2 \lambda \sqrt{\Omega}},\quad
  \beta = \frac{+ 2 C}{2 \lambda \sqrt{\Omega}},\quad
  \gamma = \frac{- 2 A}{2 \lambda \sqrt{\Omega}},\quad
  \delta = \frac{- \lambda \sqrt{\Omega} - B}{2 \lambda \sqrt{\Omega}},
\end{gather*}
(values which give $\alpha + \delta = -1$, $\alpha \delta - \beta \gamma = -1$, and therefore also
\[
  \alpha^{2} + \alpha \delta + \delta^{2} + \beta \gamma = 0,
\]
which is the condition in order that the function
$\phi x = \dfrac{\alpha x + \beta}{\gamma x + \delta}$ may be periodic of the third order,
$\phi^{3} x = x$), then, $u$ being a root of the equation, say $(a,\,b,\,c,\,d \mid u,\,1)^{3} = 0$, the
other two roots are
\[
  \phi u = \frac{\alpha u + \beta}{\gamma u + \delta}\,,
\]
and
\[
  \phi^{2} u = \phi^{-1} u
   = \frac{(\alpha^{2} + \beta \gamma)\,u + \beta(\alpha + \delta)}{\gamma(\alpha + \delta)\,u + \delta^{2} + \beta \gamma}
   = \frac{- \delta u - \beta}{- \gamma u + \alpha}\,,
\]
where observe that, by the change of $\sqrt{\Omega}$ into $-\sqrt{\Omega}$,
$\alpha,\,\beta,\,\gamma,\,\delta$ become $\delta,\,-\beta,\,-\gamma,\,\alpha$;
so that the last-mentioned value $\phi^{-1} u$ is, in fact, the value obtained from $\phi u$ by the
mere change of sign of the radical.

\hfill $52{-}2$

% ----- PDF p436 (printed 412) -----
It is to be observed that, if we have between two roots $u,\,v$ of the equation
\[
  (a,\,b,\,c,\,d \mid x,\,1)^{3} = 0,
\]
a relation $v = \dfrac{\alpha u + \beta}{\gamma u + \delta}$, where $\alpha,\,\beta,\,\gamma,\,\delta$ have given values, this implies in the first place
a relation between $a,\,b,\,c,\,d$ (and the given values of $\alpha,\,\beta,\,\gamma,\,\delta$), and it implies more-
over that $u$, and consequently also $v$, are not any roots whatever, but two determinate
roots of the equation; \emph{viz.} $u,\,v$ will be each of them expressible rationally in terms
of $a,\,b,\,c,\,d$ and $\alpha,\,\beta,\,\gamma,\,\delta$. And if, in order that $(a,\,b,\,c,\,d)$ may remain arbitrary,
we consider $\alpha,\,\beta,\,\gamma,\,\delta$ as given quantities satisfying the relation which exists between
these quantities and $(a,\,b,\,c,\,d)$, then in general we still have $u,\,v$ determinate roots
of the cubic equation. But in the foregoing solution $u$ is any root whatever of the
cubic equation.

\medskip
To examine how this is, starting from the equations
\[
  (a,\,b,\,c,\,d \mid u,\,1)^{3} = 0,\quad
  (a,\,b,\,c,\,d \mid v,\,1)^{3} = 0,\quad
  v = \frac{\alpha u + \beta}{\gamma u + \delta}\,,
\]
we have
\[
  a(u^{3} - v^{3}) + 3 b(u^{2} - v^{2}) + 3 c(u - v) = 0,
\]
and therefore
\[
  a(u^{2} + u v + v^{2}) + 3 b(u + v) + 3 c = 0,
\]
that is,
\[
  a v^{2} + (a u + 3 b)\,v + a u^{2} + 3 b u + 3 c = 0;
\]
or, writing herein for $v$ its value,
\[
  a(\alpha u + \beta)^{2} + (\alpha u + \beta)(\gamma u + \delta)(a u + 3 b) + (\gamma u + \delta)^{2}\,(a u^{2} + 3 b u + 3 c) = 0;
\]
that is,
\begin{align*}
  & a(\alpha u + \beta)^{2}
    + \alpha \gamma\,(a u^{2} + 3 b u^{2})
    + \gamma^{2}(a u^{4} + 3 b u^{3} + 3 c u^{2})\\
  & {} + (\alpha \delta + \beta \gamma)(a u^{2} + 3 b u)
    + 2 \gamma \delta\,(a u^{3} + 3 b u^{2} + 3 c u)\\
  & {} + \beta \delta\,(a u + 3 b)
    + \delta^{2}(a u^{2} + 3 b u + 3 c) = 0;
\end{align*}
or, reducing by the equation $a u^{3} + 3 b u^{2} + 3 c u + d = 0$, this is
\begin{align*}
  & a(\alpha u + \beta)^{2}
    + \alpha \gamma\,(-3 c u - d)
    + \gamma^{2}(-d u)\\
  & {} + (\alpha \delta + \beta \gamma)(a u^{2} + 3 b u)
    + \beta \gamma \delta\,(-d)\\
  & {} + \beta \delta\,(a u + 3 b)
    + \delta^{2}(a u^{2} + 3 b u + 3 c) = 0,
\end{align*}
and, collecting the terms, this is
\begin{align*}
  & u^{2}\,a\,(\alpha^{2} + \alpha \delta + \delta^{2} + \beta \gamma)\\
  & {} + u\,\{a(2 \alpha \beta + \beta \delta) + 3 b(\alpha \delta + \delta^{2} + \beta \gamma) - 3 c \alpha \gamma - d \gamma^{2}\}\\
  & {} + a \beta^{2} + 3 b \beta \delta + 3 c \delta^{2} + d(- \alpha \gamma - 2 \gamma \delta) = 0.
\end{align*}

\medskip
We can, from this equation, and the equation $a u^{3} + 3 b u^{2} + 3 c u + d = 0$, eliminate $u$, thus
obtaining a relation between $a,\,b,\,c,\,d,\,\alpha,\,\beta,\,\gamma,\,\delta$; and, this relation being satisfied, the
two equations then determine $u$ rationally in terms of these quantities.

% ----- PDF p437 (printed 413) -----
We may without loss of generality assume $\alpha \delta - \beta \gamma = 1$; and, this being so, if we
then further assume $\alpha + \delta = -1$, then we have
\[
  \alpha^{2} + \alpha \delta + \delta^{2} + \beta \gamma = 0,
\]
which is, as above appearing, the condition for $\phi^{3} x = 0$. The foregoing equation in $u$
thus becomes
\begin{align*}
  & u\,\{a \beta^{2}(\alpha - 1) - 3 b \alpha^{2} - 3 c \alpha \gamma - d \gamma^{2}\}\\
  & {} + \{a \beta^{2} + 3 b \beta \delta + 3 c \delta^{2} - d \gamma(\delta + 1)\} = 0;
\end{align*}
a linear equation giving (in a simplified form) the like results to those given by the
quadric equation; \emph{viz.} substituting in the cubic equation the value of $u$ given by the
linear equation, we have a relation between $a,\,b,\,c,\,d,\,\alpha,\,\beta,\,\gamma,\,\delta$; and, this relation
being satisfied, $u$ has the determinate value given by the linear equation.

\medskip
The only way in which $u$ can cease to have this determinate value, and so be
capable of being any root whatever of the cubic equation, is when the linear equation
becomes $0 = 0$; \emph{viz.} if
\begin{align*}
  a \beta^{2}(\alpha - 1) - 3 b \alpha^{2} - 3 c \alpha \gamma - d \gamma^{2} &= 0,\\
  a \beta^{2} + 3 b \beta \delta + 3 c \delta^{2} - d \gamma(\delta + 1) &= 0,
\end{align*}
equations which are, in fact, satisfied by the foregoing values of $\alpha,\,\beta,\,\gamma,\,\delta$, as may be
verified without difficulty.

\medskip
It is to be remarked that if, instead of the root $u$ and the equation
$v = \dfrac{\alpha u + \beta}{\gamma u + \delta}$, we consider the root $v$ and the equation
$u = \dfrac{- \delta v - \beta}{- \gamma v + \alpha}$, then, instead of $\alpha,\,\beta,\,\gamma,\,\delta$, we
have $\delta,\,-\beta,\,-\gamma,\,\alpha$, and the corresponding equations are
\begin{align*}
  a \beta^{2}(\delta - 1) - 3 b \delta^{2} + 3 c \gamma \delta + d \gamma^{2} &= 0,\\
  a \beta^{2} + 3 b \alpha \beta + 3 c \alpha^{2} + d \gamma(\alpha - 1) &= 0,
\end{align*}
equations which are also satisfied by the foregoing values of $\alpha,\,\beta,\,\gamma,\,\delta$. And the four
equations, together with $\alpha \delta - \beta \gamma = 1$ and $\alpha + \delta = -1$, are more than sufficient to determine
the foregoing values of $\alpha,\,\beta,\,\gamma,\,\delta$.

\medskip
But we further verify without difficulty that the foregoing values of $\alpha,\,\beta,\,\gamma,\,\delta$ give
identically
\[
  (a,\,b,\,c,\,d \mid \alpha x + \beta y,\,\gamma x + \delta y)^{3} = -(a,\,b,\,c,\,d \mid x,\,y)^{3},
\]
or the formul\ae\ lead to an automorphic transformation of the binary cubic
$(a,\,b,\,c,\,d \mid x,\,y)^{3}$.
And conversely, starting from the notion of the automorphic transformation of the binary
cubic, we ought to be able to obtain the foregoing formul\ae.

\medskip
For greater convenience, I write the equation of transformation in the form
\[
  (a,\,b,\,c,\,d \mid \alpha x + \beta y,\,\gamma x + \delta y)^{3} = -\,\theta\,(a,\,b,\,c,\,d \mid x,\,y)^{3};
\]

\end{document}
